%! Observational Studies — Unit 1, Lesson 1.5 — Homework
% PILOT SHAPE (2026-09-06): modeled on AP Statistics lesson 1.4 minus its AP Practice page.
% This is the lesson's SCORED individual practice and the LAST component in the packet.
% Students START IT IN CLASS in the last ten minutes of the period, alone, while the
% teacher circulates for the second formative read; they finish it at home. Packet pages are
% the default; a Desmos activity may replace them on a given day (decided at Close & Assign).
% THIRD CONTEXT: the notes teach on the Harbor Point pool, the Guided Practice runs on
% Riverbend's Study Center, and these items are on the Millbrook Public Library — recall is
% not the skill.
% TWO PAGES MAX (user decision 2026-09-06): two boxes — items 1-2 on page 1, items 3-6 and
% the spiralbox on page 2 — so no item breaks across a page.
% Item 4's two rows are the deliberate CONTRAST PAIR (the association sentence, supported;
% the newsletter's causal sentence, not). Item 5 is the arrow backwards, from the previous
% summer's records. Item 6 is the crux head-on: no bias at all, and still no causation.
% Item 3 is the spiral to 1.2 (pool the sample, scale to the population, call it an estimate).
% THE DATA — Millbrook Public Library: 1,200 card holders; random phone sample of 240, all
%   reached; 90 signed up for the Summer Reading Program, 72 read 10+ books -> 72/90 = 80%;
%   150 did not sign up, 48 read 10+ -> 48/150 = 32%; gap 48 points; pooled 120/240 = 50%;
%   0.50 x 1,200 = 600; previous summer 63/90 = 70% vs 30/150 = 20%. Verified in Python.
% PAGE LOCKSTEP: the key is generated from this file; every work block is byte-identical.
\documentclass[12pt]{article}
\usepackage{saar-article}
\usepackage{saar-boxes}
\newcommand{\TallMath}[1]{$\displaystyle #1\rule[-1.4em]{0pt}{3.2em}$}
% Word bank strip for every fill-in cluster (user request 2026-09-06). Defined per document;
% the key inherits it and prints the same strip, so heights match.
\newcommand{\wordbank}[1]{\par\vspace{2pt}\noindent\fcolorbox{goldacc}{goldbg}{\parbox{\dimexpr\linewidth-2\fboxsep-2\fboxrule\relax}{\small\textbf{Word bank:} #1}}\par\vspace{4pt}}

\begin{document}

\pageheader{Unit 1, Lesson 1.5}{Homework}
% NAMESTRIP: no \namedateperiod here. The student writes their name once, on the
% cover this component is stapled behind.

\vspace{0.12in}

% Authored identically in homework/ and homework_key/.
\begin{remindbox}
\small
\textbf{This is your graded homework.} It is scored out of 12 and due at the start of the
first class after two study halls. You will start it in the last minutes of the period, so ask about anything that stalls
you before you leave, then finish it on your own.
\end{remindbox}

\vspace{0.08in}

\boxguard[20]
\begin{scenariobox}[Millbrook Public Library --- Does the Summer Reading Program Work?]{cerulean}
Millbrook Public Library has $1{,}200$ card holders. Every summer it runs a free Summer
Reading Program that any card holder may sign up for. In September the director draws a
random sample of $240$ card holders from the full membership list and reaches every one of
them by phone.

\vspace{2pt}
Of the $240$, $90$ had signed up for the program and $150$ had not. Of the $90$ who signed
up, $72$ read at least $10$ books over the summer. Of the $150$ who did not, $48$ did.
\end{scenariobox}

\vspace{0.08in}

\boxguard
\begin{notesbox}{Practice}
Every answer names the group it is about --- the $90$ who signed up, the $150$ who did not,
or all $1{,}200$ card holders.
\wordbank{observational \;\textbullet\; the card holders themselves \;\textbullet\; signed up for the program \;\textbullet\; read at least 10 books \;\textbullet\; 48 \;\textbullet\; the program group}
\begin{enumerate}[leftmargin=1.6em, itemsep=9pt, topsep=3pt]

  \item The director did not put anybody into the program --- card holders signed themselves up.

        \par\vspace{4pt}
        What kind of study is this? \blank{3.4cm} \hfill Who decided the groups? \blank{4.6cm}
        \par\vspace{8pt}
        Explanatory variable: \blank{5.2cm} \hfill Response variable: \blank{5cm}

  \item Compute the reading rate for each group --- watch the two denominators --- then the
        gap.

        \begin{work}
          \dfrac{72}{90} &= 0.80 = 80\% \text{ of the program group} \\
          \dfrac{48}{150} &= 0.32 = 32\% \text{ of the others}
        \end{work}

        \par\vspace{2pt}
        The \blank{3.2cm} group is higher by \blank{1.8cm} percentage points.

\end{enumerate}
\end{notesbox}

\boxguard[30]
\begin{notesbox}{Practice, continued}
\begin{enumerate}[leftmargin=1.6em, itemsep=5pt, topsep=2pt, start=3]

  \item Now use the whole sample. How many of the $240$ read at least $10$ books, what
        percent of the sample is that, and about how many of all $1{,}200$ card holders does
        that predict?

        \begin{work}
          \dfrac{72 + 48}{240} &= \dfrac{120}{240} = 0.50 = 50\% \\
          0.50 \times 1200 &= 600 \text{ card holders --- an estimate}
        \end{work}

  \item Same data, two sentences. Decide which one the study supports, and say why.

        \begin{center}
        \renewcommand{\arraystretch}{1.3}
        \begin{tabularx}{\linewidth}{@{} X p{2.2cm} >{\raggedright\arraybackslash}p{4.6cm} @{}}
        \toprule
        \textbf{Statement} & \textbf{Supported?} & \textbf{Why / why not} \\
        \midrule
        Card holders who signed up read more books than card holders who did not.
          & \blank{1.8cm} & \blank{4.2cm} \\
        The Summer Reading Program more than doubles how much our members read.
          & \blank{1.8cm} & \blank{4.2cm} \\
        \bottomrule
        \end{tabularx}
        \end{center}

  \item The director checks one more record: how much each card holder read the
        \emph{previous} summer, before the program existed. Of the $90$ who signed up, $63$
        had already read $10$ or more books. Of the $150$ who did not, $30$ had.

        \par\vspace{4pt}
        Signed up: \blank{3.6cm} \qquad Did not sign up: \blank{3.6cm}

        \par\vspace{6pt}
        Which of the three stories is this --- and does it make the program look
        \emph{better} or \emph{worse} than it really is?
        \par\writelines{2}

  \item The sample was drawn at random from the full list, every person answered, and every
        count is correct --- there is \emph{no bias} in this study. Explain why it
        \emph{still} cannot show that the program caused the extra reading.
        \par\writelines{2}

\end{enumerate}
\end{notesbox}

\boxguard[5]   % the two-line spiralbox needs about five baselines; a larger guard pushes it to a third page
\begin{spiralbox}
\textbf{Next lesson: Experimental Design.} An observational study cannot reach causation
because the groups build themselves. Next class we build them instead --- treatment and
control, \emph{assigned} at random --- and meet the four principles that make an experiment
worth trusting.
\end{spiralbox}

\end{document}
